\section{Learning Objectives}

After completing this chapter, readers will be able to:
\begin{itemize}
    \item Design and implement end-to-end ML systems for real-world applications
    \item Understand architecture patterns across different industries
    \item Apply domain-specific best practices and optimization techniques
    \item Learn from production challenges and their solutions
    \item Recognize common patterns and anti-patterns in ML system design
    \item Implement scalable data pipelines for various use cases
\end{itemize}

\section{Introduction}

This chapter presents four comprehensive industry case studies that demonstrate end-to-end ML system implementations. Each case study covers the complete lifecycle from requirements gathering through production deployment, including real-world challenges and solutions.

\textbf{Case Studies Covered:}

\begin{enumerate}
    \item \textbf{E-Commerce Recommendation System:} Personalized product recommendations at scale
    \item \textbf{Financial Fraud Detection:} Real-time transaction monitoring and fraud prevention
    \item \textbf{Healthcare Diagnostic System:} Medical image analysis for disease detection
    \item \textbf{Autonomous Vehicle Perception:} Multi-sensor fusion for object detection
\end{enumerate}

Each case study follows a consistent structure: business context, requirements, architecture design, implementation, deployment, monitoring, results, and lessons learned.

\section{Case Study 1: E-Commerce Recommendation System}

\subsection{Business Context and Requirements}

\textbf{Company:} Mid-sized e-commerce platform with 5 million monthly active users\\
\textbf{Challenge:} Low conversion rates and poor product discovery\\
\textbf{Goal:} Increase revenue through personalized product recommendations

\textbf{Key Requirements:}

\begin{itemize}
    \item \textbf{Latency:} Recommendations must load in <100ms (p99)
    \item \textbf{Scale:} Handle 10,000 requests per second during peak hours
    \item \textbf{Personalization:} Adapt to user behavior in real-time
    \item \textbf{Business Constraints:} Incorporate inventory, margins, and business rules
    \item \textbf{Cold Start:} Handle new users and new products effectively
    \item \textbf{Diversity:} Avoid filter bubbles, show varied recommendations
\end{itemize}

\subsection{System Architecture}

\begin{figure}[h]
\centering
\resizebox{\textwidth}{!}{%
\begin{tikzpicture}[
    node distance=1.5cm,
    component/.style={rectangle, draw, fill=blue!20, text width=3cm, text centered, minimum height=1cm},
    storage/.style={cylinder, draw, fill=green!20, text width=2.5cm, text centered, minimum height=1cm, aspect=0.3},
    process/.style={rectangle, draw, fill=yellow!20, text width=3cm, text centered, minimum height=0.8cm}
]

% Data sources
\node[storage] (clicks) {User Events\\(Kafka)};
\node[storage, right=of clicks] (catalog) {Product\\Catalog};
\node[storage, right=of catalog] (orders) {Order\\History};

% Processing
\node[process, below=of clicks] (stream) {Stream\\Processing};
\node[process, below=of catalog] (batch) {Batch\\Training};
\node[storage, below=of stream] (features) {Feature\\Store};

% Models
\node[component, below=of features] (cf) {Collaborative\\Filtering};
\node[component, right=of cf] (content) {Content-Based};
\node[component, left=of cf] (deep) {Deep\\Learning};

% Serving
\node[component, below=of cf] (ensemble) {Ensemble\\Ranker};
\node[component, below=of ensemble] (api) {Serving\\API};

% Arrows
\draw[->] (clicks) -- (stream);
\draw[->] (catalog) -- (batch);
\draw[->] (orders) -- (batch);
\draw[->] (stream) -- (features);
\draw[->] (batch) -- (features);
\draw[->] (features) -- (cf);
\draw[->] (features) -- (content);
\draw[->] (features) -- (deep);
\draw[->] (cf) -- (ensemble);
\draw[->] (content) -- (ensemble);
\draw[->] (deep) -- (ensemble);
\draw[->] (ensemble) -- (api);

\end{tikzpicture}}
\caption{E-Commerce Recommendation System Architecture}
\end{figure}

\subsection{Data Pipeline Implementation}

\begin{lstlisting}[language=Python, caption=Event Streaming and Feature Engineering]
#!/usr/bin/env python3
"""
E-Commerce Recommendation System - Data Pipeline
Real-time and batch feature engineering for personalized recommendations
"""

from kafka import KafkaConsumer, KafkaProducer
import json
from typing import Dict, List, Optional
from datetime import datetime, timedelta
from dataclasses import dataclass, field
import pandas as pd
import numpy as np
from redis import Redis
from collections import defaultdict
import logging

@dataclass
class UserEvent:
    """User interaction event"""
    user_id: str
    event_type: str  # view, click, add_to_cart, purchase
    product_id: str
    timestamp: datetime
    session_id: str
    metadata: Dict = field(default_factory=dict)

@dataclass
class Product:
    """Product information"""
    product_id: str
    category: str
    subcategory: str
    price: float
    brand: str
    in_stock: bool
    attributes: Dict = field(default_factory=dict)

@dataclass
class UserFeatures:
    """Aggregated user features"""
    user_id: str
    favorite_categories: List[str]
    avg_price_point: float
    preferred_brands: List[str]
    recent_views: List[str]
    recent_purchases: List[str]
    purchase_frequency: float
    last_active: datetime

class RecommendationDataPipeline:
    """Data pipeline for recommendation system"""

    def __init__(
        self,
        kafka_bootstrap_servers: List[str],
        redis_host: str,
        redis_port: int = 6379
    ):
        self.kafka_consumer = KafkaConsumer(
            'user-events',
            bootstrap_servers=kafka_bootstrap_servers,
            value_deserializer=lambda m: json.loads(m.decode('utf-8')),
            group_id='recommendation-pipeline'
        )

        self.kafka_producer = KafkaProducer(
            bootstrap_servers=kafka_bootstrap_servers,
            value_serializer=lambda m: json.dumps(m).encode('utf-8')
        )

        self.redis_client = Redis(host=redis_host, port=redis_port, decode_responses=True)
        self.logger = logging.getLogger(__name__)

    def process_event_stream(self):
        """Process real-time user events"""

        for message in self.kafka_consumer:
            try:
                event_data = message.value
                event = UserEvent(
                    user_id=event_data['user_id'],
                    event_type=event_data['event_type'],
                    product_id=event_data['product_id'],
                    timestamp=datetime.fromisoformat(event_data['timestamp']),
                    session_id=event_data['session_id'],
                    metadata=event_data.get('metadata', {})
                )

                # Update real-time features
                self._update_user_features(event)

                # Update product popularity
                self._update_product_stats(event)

                # Generate real-time signals
                if event.event_type == 'purchase':
                    self._trigger_similar_products(event)

            except Exception as e:
                self.logger.error(f"Error processing event: {e}")

    def _update_user_features(self, event: UserEvent):
        """Update real-time user features in Redis"""

        user_key = f"user:{event.user_id}"

        # Update recent views (sliding window)
        if event.event_type in ['view', 'click']:
            self.redis_client.lpush(f"{user_key}:recent_views", event.product_id)
            self.redis_client.ltrim(f"{user_key}:recent_views", 0, 49)  # Keep last 50

        # Update purchases
        if event.event_type == 'purchase':
            self.redis_client.lpush(f"{user_key}:purchases", event.product_id)
            self.redis_client.lpush(f"{user_key}:purchase_times", event.timestamp.isoformat())

        # Update category preferences
        product_category = self._get_product_category(event.product_id)
        if product_category:
            self.redis_client.zincrby(
                f"{user_key}:categories",
                1,
                product_category
            )

        # Update session activity
        self.redis_client.hset(
            f"{user_key}:session:{event.session_id}",
            mapping={
                'last_activity': event.timestamp.isoformat(),
                'event_count': self.redis_client.hincrby(
                    f"{user_key}:session:{event.session_id}",
                    'event_count',
                    1
                )
            }
        )
        self.redis_client.expire(f"{user_key}:session:{event.session_id}", 3600)

    def _update_product_stats(self, event: UserEvent):
        """Update product statistics for trending items"""

        product_key = f"product:{event.product_id}"

        # Increment view/purchase counters with decay
        current_hour = datetime.now().strftime("%Y-%m-%d-%H")

        if event.event_type == 'view':
            self.redis_client.hincrby(f"{product_key}:stats:{current_hour}", 'views', 1)
        elif event.event_type == 'purchase':
            self.redis_client.hincrby(f"{product_key}:stats:{current_hour}", 'purchases', 1)

        # Set expiry on hourly stats
        self.redis_client.expire(f"{product_key}:stats:{current_hour}", 86400 * 7)

        # Update trending score (combines recency and popularity)
        trending_score = self._calculate_trending_score(event.product_id)
        self.redis_client.zadd('trending_products', {event.product_id: trending_score})

    def _calculate_trending_score(self, product_id: str) -> float:
        """Calculate trending score based on recent activity"""

        product_key = f"product:{product_id}"
        score = 0.0

        # Look at last 24 hours
        for hours_ago in range(24):
            timestamp = (datetime.now() - timedelta(hours=hours_ago)).strftime("%Y-%m-%d-%H")
            stats = self.redis_client.hgetall(f"{product_key}:stats:{timestamp}")

            if stats:
                views = int(stats.get('views', 0))
                purchases = int(stats.get('purchases', 0))

                # Decay factor: recent hours weighted more
                decay = 0.95 ** hours_ago

                # Weighted score
                score += (views * 1.0 + purchases * 10.0) * decay

        return score

    def _trigger_similar_products(self, event: UserEvent):
        """Trigger batch job to update similar products after purchase"""

        message = {
            'type': 'purchase',
            'user_id': event.user_id,
            'product_id': event.product_id,
            'timestamp': event.timestamp.isoformat()
        }

        self.kafka_producer.send('recommendation-triggers', value=message)

    def _get_product_category(self, product_id: str) -> Optional[str]:
        """Get product category from cache or database"""

        category = self.redis_client.get(f"product:{product_id}:category")
        if category:
            return category

        # Fallback: fetch from database (not shown)
        return None

    def get_user_features(self, user_id: str) -> UserFeatures:
        """Get aggregated user features for recommendations"""

        user_key = f"user:{user_id}"

        # Get recent views
        recent_views = self.redis_client.lrange(f"{user_key}:recent_views", 0, 19)

        # Get recent purchases
        recent_purchases = self.redis_client.lrange(f"{user_key}:purchases", 0, 9)

        # Get favorite categories
        categories = self.redis_client.zrevrange(
            f"{user_key}:categories",
            0, 4,
            withscores=True
        )
        favorite_categories = [cat for cat, score in categories]

        # Calculate purchase frequency
        purchase_times = self.redis_client.lrange(f"{user_key}:purchase_times", 0, -1)
        purchase_frequency = self._calculate_purchase_frequency(purchase_times)

        # Get average price point
        avg_price = self._calculate_avg_price_point(recent_purchases)

        return UserFeatures(
            user_id=user_id,
            favorite_categories=favorite_categories,
            avg_price_point=avg_price,
            preferred_brands=[],  # Computed similarly
            recent_views=recent_views,
            recent_purchases=recent_purchases,
            purchase_frequency=purchase_frequency,
            last_active=datetime.now()
        )

    def _calculate_purchase_frequency(self, purchase_times: List[str]) -> float:
        """Calculate purchases per week"""

        if len(purchase_times) < 2:
            return 0.0

        timestamps = [datetime.fromisoformat(t) for t in purchase_times[:10]]
        if len(timestamps) < 2:
            return 0.0

        time_span = (timestamps[0] - timestamps[-1]).total_seconds() / (7 * 86400)
        if time_span == 0:
            return 0.0

        return len(timestamps) / time_span

    def _calculate_avg_price_point(self, product_ids: List[str]) -> float:
        """Calculate user's average price point"""

        prices = []
        for product_id in product_ids[:20]:
            price = self.redis_client.get(f"product:{product_id}:price")
            if price:
                prices.append(float(price))

        return np.mean(prices) if prices else 0.0

class BatchFeatureEngineering:
    """Batch feature engineering for model training"""

    def __init__(self, data_warehouse_conn):
        self.conn = data_warehouse_conn

    def create_training_dataset(
        self,
        start_date: datetime,
        end_date: datetime
    ) -> pd.DataFrame:
        """Create training dataset from historical data"""

        # User features
        user_features = self._compute_user_features(start_date, end_date)

        # Product features
        product_features = self._compute_product_features(start_date, end_date)

        # Interaction matrix
        interactions = self._build_interaction_matrix(start_date, end_date)

        # Combine features
        training_data = self._combine_features(
            user_features,
            product_features,
            interactions
        )

        return training_data

    def _compute_user_features(
        self,
        start_date: datetime,
        end_date: datetime
    ) -> pd.DataFrame:
        """Compute aggregated user features"""

        query = """
        SELECT
            user_id,
            COUNT(DISTINCT session_id) as session_count,
            COUNT(*) as total_events,
            COUNT(DISTINCT product_id) as unique_products_viewed,
            COUNT(CASE WHEN event_type = 'purchase' THEN 1 END) as purchase_count,
            AVG(CASE WHEN event_type = 'purchase' THEN product_price END) as avg_purchase_price,
            MODE() WITHIN GROUP (ORDER BY product_category) as favorite_category,
            MAX(timestamp) as last_active
        FROM user_events
        WHERE timestamp BETWEEN %s AND %s
        GROUP BY user_id
        """

        return pd.read_sql(query, self.conn, params=[start_date, end_date])

    def _compute_product_features(
        self,
        start_date: datetime,
        end_date: datetime
    ) -> pd.DataFrame:
        """Compute product features"""

        query = """
        SELECT
            p.product_id,
            p.category,
            p.subcategory,
            p.price,
            p.brand,
            COUNT(DISTINCT ue.user_id) as unique_viewers,
            COUNT(*) as total_views,
            COUNT(CASE WHEN ue.event_type = 'purchase' THEN 1 END) as purchase_count,
            COUNT(CASE WHEN ue.event_type = 'purchase' THEN 1 END)::FLOAT /
                NULLIF(COUNT(*), 0) as conversion_rate
        FROM products p
        LEFT JOIN user_events ue ON p.product_id = ue.product_id
            AND ue.timestamp BETWEEN %s AND %s
        GROUP BY p.product_id, p.category, p.subcategory, p.price, p.brand
        """

        return pd.read_sql(query, self.conn, params=[start_date, end_date])

    def _build_interaction_matrix(
        self,
        start_date: datetime,
        end_date: datetime
    ) -> pd.DataFrame:
        """Build user-product interaction matrix"""

        query = """
        SELECT
            user_id,
            product_id,
            SUM(CASE WHEN event_type = 'view' THEN 1 ELSE 0 END) as views,
            SUM(CASE WHEN event_type = 'click' THEN 1 ELSE 0 END) as clicks,
            SUM(CASE WHEN event_type = 'add_to_cart' THEN 1 ELSE 0 END) as cart_adds,
            SUM(CASE WHEN event_type = 'purchase' THEN 1 ELSE 0 END) as purchases,
            MAX(timestamp) as last_interaction
        FROM user_events
        WHERE timestamp BETWEEN %s AND %s
        GROUP BY user_id, product_id
        """

        return pd.read_sql(query, self.conn, params=[start_date, end_date])

    def _combine_features(
        self,
        user_features: pd.DataFrame,
        product_features: pd.DataFrame,
        interactions: pd.DataFrame
    ) -> pd.DataFrame:
        """Combine all features into training dataset"""

        # Merge user and product features with interactions
        training_data = interactions.merge(
            user_features,
            on='user_id',
            how='left'
        ).merge(
            product_features,
            on='product_id',
            how='left'
        )

        # Create target variable (implicit feedback)
        training_data['target'] = (
            training_data['views'] * 0.1 +
            training_data['clicks'] * 0.3 +
            training_data['cart_adds'] * 0.5 +
            training_data['purchases'] * 1.0
        )

        # Create negative samples
        training_data = self._add_negative_samples(training_data)

        return training_data

    def _add_negative_samples(self, positive_samples: pd.DataFrame) -> pd.DataFrame:
        """Add negative samples for training"""

        # For each user, sample products they didn't interact with
        # This is simplified - in production, use more sophisticated negative sampling

        all_users = positive_samples['user_id'].unique()
        all_products = positive_samples['product_id'].unique()

        negative_samples = []

        for user_id in all_users:
            user_products = set(positive_samples[
                positive_samples['user_id'] == user_id
            ]['product_id'])

            # Sample 4 negative products per positive
            num_negatives = len(user_products) * 4
            negative_products = np.random.choice(
                list(set(all_products) - user_products),
                size=min(num_negatives, len(all_products) - len(user_products)),
                replace=False
            )

            for product_id in negative_products:
                negative_samples.append({
                    'user_id': user_id,
                    'product_id': product_id,
                    'target': 0
                })

        negative_df = pd.DataFrame(negative_samples)

        return pd.concat([positive_samples, negative_df], ignore_index=True)

# Example usage
if __name__ == "__main__":
    pipeline = RecommendationDataPipeline(
        kafka_bootstrap_servers=['localhost:9092'],
        redis_host='localhost'
    )

    # Start processing events
    # pipeline.process_event_stream()

    # Get user features
    features = pipeline.get_user_features('user_12345')
    print(f"User categories: {features.favorite_categories}")
    print(f"Recent views: {len(features.recent_views)}")
    print(f"Purchase frequency: {features.purchase_frequency:.2f} per week")
\end{lstlisting}

\subsection{Model Implementation}

\begin{lstlisting}[language=Python, caption=Hybrid Recommendation Models]
#!/usr/bin/env python3
"""
E-Commerce Recommendation Models
Hybrid approach combining collaborative filtering, content-based, and deep learning
"""

import torch
import torch.nn as nn
import numpy as np
from typing import List, Dict, Tuple
from sklearn.metrics.pairwise import cosine_similarity
import pandas as pd
from scipy.sparse import csr_matrix
from implicit.als import AlternatingLeastSquares
from implicit.bpr import BayesianPersonalizedRanking

class CollaborativeFilteringModel:
    """Matrix factorization for collaborative filtering"""

    def __init__(self, factors: int = 64, regularization: float = 0.01):
        self.factors = factors
        self.model = AlternatingLeastSquares(
            factors=factors,
            regularization=regularization,
            iterations=15
        )
        self.user_id_map = {}
        self.product_id_map = {}

    def fit(self, interactions: pd.DataFrame):
        """Train collaborative filtering model"""

        # Create user and product mappings
        unique_users = interactions['user_id'].unique()
        unique_products = interactions['product_id'].unique()

        self.user_id_map = {uid: idx for idx, uid in enumerate(unique_users)}
        self.product_id_map = {pid: idx for idx, pid in enumerate(unique_products)}

        # Create sparse interaction matrix
        user_indices = interactions['user_id'].map(self.user_id_map)
        product_indices = interactions['product_id'].map(self.product_id_map)
        values = interactions['target'].values

        interaction_matrix = csr_matrix(
            (values, (user_indices, product_indices)),
            shape=(len(unique_users), len(unique_products))
        )

        # Train model
        self.model.fit(interaction_matrix)

    def recommend(
        self,
        user_id: str,
        n: int = 10,
        filter_already_liked: bool = True
    ) -> List[Tuple[str, float]]:
        """Generate recommendations for user"""

        if user_id not in self.user_id_map:
            return []

        user_idx = self.user_id_map[user_id]

        # Get recommendations
        product_indices, scores = self.model.recommend(
            user_idx,
            None,  # All items
            N=n,
            filter_already_liked_items=filter_already_liked
        )

        # Map back to product IDs
        reverse_product_map = {idx: pid for pid, idx in self.product_id_map.items()}
        recommendations = [
            (reverse_product_map[idx], float(score))
            for idx, score in zip(product_indices, scores)
        ]

        return recommendations

    def similar_items(self, product_id: str, n: int = 10) -> List[Tuple[str, float]]:
        """Find similar products"""

        if product_id not in self.product_id_map:
            return []

        product_idx = self.product_id_map[product_id]

        # Get similar items
        similar_indices, scores = self.model.similar_items(product_idx, N=n+1)

        # Exclude the item itself
        similar_indices = similar_indices[1:]
        scores = scores[1:]

        reverse_product_map = {idx: pid for pid, idx in self.product_id_map.items()}
        similar_products = [
            (reverse_product_map[idx], float(score))
            for idx, score in zip(similar_indices, scores)
        ]

        return similar_products

class ContentBasedModel:
    """Content-based filtering using product features"""

    def __init__(self):
        self.product_vectors = {}
        self.product_features = None

    def fit(self, product_features: pd.DataFrame):
        """Train content-based model"""

        self.product_features = product_features.set_index('product_id')

        # Create product vectors from features
        for product_id, row in self.product_features.iterrows():
            vector = self._create_product_vector(row)
            self.product_vectors[product_id] = vector

    def _create_product_vector(self, product_row) -> np.ndarray:
        """Create feature vector for product"""

        # Combine categorical and numerical features
        features = []

        # One-hot encode category
        category_features = self._one_hot_encode(
            product_row['category'],
            self.product_features['category'].unique()
        )
        features.extend(category_features)

        # Normalize price
        price_normalized = (product_row['price'] - self.product_features['price'].mean()) / \
                          self.product_features['price'].std()
        features.append(price_normalized)

        # Brand encoding
        brand_features = self._one_hot_encode(
            product_row['brand'],
            self.product_features['brand'].unique()
        )
        features.extend(brand_features)

        return np.array(features)

    def _one_hot_encode(self, value, all_values):
        """One-hot encode a categorical value"""
        return [1 if v == value else 0 for v in all_values]

    def recommend(
        self,
        user_profile: Dict,
        n: int = 10,
        exclude_products: List[str] = None
    ) -> List[Tuple[str, float]]:
        """Recommend products based on user profile"""

        if exclude_products is None:
            exclude_products = []

        # Create user vector from profile
        user_vector = self._create_user_vector(user_profile)

        # Calculate similarity to all products
        similarities = {}
        for product_id, product_vector in self.product_vectors.items():
            if product_id in exclude_products:
                continue

            similarity = cosine_similarity(
                user_vector.reshape(1, -1),
                product_vector.reshape(1, -1)
            )[0][0]

            similarities[product_id] = similarity

        # Sort by similarity
        recommendations = sorted(
            similarities.items(),
            key=lambda x: x[1],
            reverse=True
        )[:n]

        return recommendations

    def _create_user_vector(self, user_profile: Dict) -> np.ndarray:
        """Create user vector from their interaction history"""

        # Average of viewed/purchased product vectors
        product_ids = user_profile.get('recent_views', []) + \
                     user_profile.get('recent_purchases', [])

        if not product_ids:
            # Return zero vector for cold start
            sample_vector = next(iter(self.product_vectors.values()))
            return np.zeros_like(sample_vector)

        vectors = [
            self.product_vectors[pid]
            for pid in product_ids
            if pid in self.product_vectors
        ]

        if not vectors:
            sample_vector = next(iter(self.product_vectors.values()))
            return np.zeros_like(sample_vector)

        return np.mean(vectors, axis=0)

    def similar_items(self, product_id: str, n: int = 10) -> List[Tuple[str, float]]:
        """Find similar products by content"""

        if product_id not in self.product_vectors:
            return []

        product_vector = self.product_vectors[product_id]

        similarities = {}
        for pid, pvector in self.product_vectors.items():
            if pid == product_id:
                continue

            similarity = cosine_similarity(
                product_vector.reshape(1, -1),
                pvector.reshape(1, -1)
            )[0][0]

            similarities[pid] = similarity

        similar_products = sorted(
            similarities.items(),
            key=lambda x: x[1],
            reverse=True
        )[:n]

        return similar_products

class DeepRecommenderModel(nn.Module):
    """Deep learning model for recommendations"""

    def __init__(
        self,
        num_users: int,
        num_products: int,
        embedding_dim: int = 64,
        hidden_dims: List[int] = [128, 64, 32]
    ):
        super().__init__()

        self.user_embedding = nn.Embedding(num_users, embedding_dim)
        self.product_embedding = nn.Embedding(num_products, embedding_dim)

        # Deep neural network
        layers = []
        input_dim = embedding_dim * 2

        for hidden_dim in hidden_dims:
            layers.append(nn.Linear(input_dim, hidden_dim))
            layers.append(nn.ReLU())
            layers.append(nn.Dropout(0.2))
            input_dim = hidden_dim

        layers.append(nn.Linear(input_dim, 1))
        layers.append(nn.Sigmoid())

        self.mlp = nn.Sequential(*layers)

    def forward(self, user_ids, product_ids):
        """Forward pass"""

        user_embeds = self.user_embedding(user_ids)
        product_embeds = self.product_embedding(product_ids)

        # Concatenate embeddings
        x = torch.cat([user_embeds, product_embeds], dim=1)

        # Pass through MLP
        output = self.mlp(x)

        return output.squeeze()

class HybridRecommender:
    """Hybrid recommender combining multiple models"""

    def __init__(self):
        self.cf_model = None
        self.content_model = None
        self.deep_model = None

        # Weights for ensemble
        self.cf_weight = 0.4
        self.content_weight = 0.3
        self.deep_weight = 0.3

    def train(
        self,
        interactions: pd.DataFrame,
        product_features: pd.DataFrame,
        user_features: pd.DataFrame
    ):
        """Train all models"""

        print("Training collaborative filtering model...")
        self.cf_model = CollaborativeFilteringModel(factors=64)
        self.cf_model.fit(interactions)

        print("Training content-based model...")
        self.content_model = ContentBasedModel()
        self.content_model.fit(product_features)

        print("Training deep learning model...")
        # Deep model training would go here
        # Omitted for brevity

        print("All models trained successfully")

    def recommend(
        self,
        user_id: str,
        user_features: Dict,
        n: int = 10,
        exclude_products: List[str] = None
    ) -> List[Tuple[str, float]]:
        """Generate hybrid recommendations"""

        if exclude_products is None:
            exclude_products = []

        # Get recommendations from each model
        cf_recs = self.cf_model.recommend(user_id, n=50, filter_already_liked=True)
        content_recs = self.content_model.recommend(user_features, n=50, exclude_products=exclude_products)

        # Combine scores
        combined_scores = defaultdict(float)

        for product_id, score in cf_recs:
            if product_id not in exclude_products:
                combined_scores[product_id] += score * self.cf_weight

        for product_id, score in content_recs:
            if product_id not in exclude_products:
                combined_scores[product_id] += score * self.content_weight

        # Sort by combined score
        recommendations = sorted(
            combined_scores.items(),
            key=lambda x: x[1],
            reverse=True
        )[:n]

        return recommendations

    def diversify_recommendations(
        self,
        recommendations: List[Tuple[str, float]],
        product_features: pd.DataFrame,
        diversity_weight: float = 0.3
    ) -> List[Tuple[str, float]]:
        """Add diversity to recommendations"""

        if not recommendations:
            return recommendations

        diversified = [recommendations[0]]  # Start with top recommendation
        remaining = recommendations[1:]

        while len(diversified) < len(recommendations) and remaining:
            best_candidate = None
            best_score = -float('inf')

            for product_id, relevance_score in remaining:
                # Calculate diversity score (dissimilarity to already selected)
                diversity_score = self._calculate_diversity(
                    product_id,
                    [p[0] for p in diversified],
                    product_features
                )

                # Combined score
                combined = (1 - diversity_weight) * relevance_score + \
                          diversity_weight * diversity_score

                if combined > best_score:
                    best_score = combined
                    best_candidate = (product_id, relevance_score)

            if best_candidate:
                diversified.append(best_candidate)
                remaining.remove(best_candidate)

        return diversified

    def _calculate_diversity(
        self,
        product_id: str,
        selected_products: List[str],
        product_features: pd.DataFrame
    ) -> float:
        """Calculate diversity score for a product"""

        if product_id not in product_features.index:
            return 0.0

        product_category = product_features.loc[product_id, 'category']

        # Calculate category diversity
        selected_categories = [
            product_features.loc[pid, 'category']
            for pid in selected_products
            if pid in product_features.index
        ]

        if not selected_categories:
            return 1.0

        # Diversity = fraction of different categories
        different_categories = sum(1 for cat in selected_categories if cat != product_category)
        diversity = different_categories / len(selected_categories)

        return diversity

# Example usage
if __name__ == "__main__":
    # Load data
    interactions = pd.read_csv('interactions.csv')
    product_features = pd.read_csv('products.csv')
    user_features = pd.read_csv('users.csv')

    # Train hybrid model
    recommender = HybridRecommender()
    recommender.train(interactions, product_features, user_features)

    # Generate recommendations
    user_id = 'user_12345'
    user_profile = {
        'recent_views': ['prod_1', 'prod_5', 'prod_7'],
        'recent_purchases': ['prod_3']
    }

    recommendations = recommender.recommend(user_id, user_profile, n=10)

    print(f"Top 10 recommendations for {user_id}:")
    for product_id, score in recommendations:
        print(f"  {product_id}: {score:.4f}")

    # Add diversity
    diverse_recs = recommender.diversify_recommendations(
        recommendations,
        product_features,
        diversity_weight=0.3
    )

    print("\nDiversified recommendations:")
    for product_id, score in diverse_recs:
        print(f"  {product_id}: {score:.4f}")
\end{lstlisting}

\subsection{Serving and Deployment}

\begin{lstlisting}[language=Python, caption=Recommendation Serving API]
#!/usr/bin/env python3
"""
Recommendation Serving API
Low-latency serving infrastructure with caching and business rules
"""

from fastapi import FastAPI, HTTPException, Query
from pydantic import BaseModel
from typing import List, Optional, Dict
import redis
import pickle
import time
from datetime import datetime
import logging

app = FastAPI(title="Recommendation API")

class RecommendationRequest(BaseModel):
    """Recommendation request"""
    user_id: str
    num_recommendations: int = 10
    context: Optional[Dict] = None
    exclude_products: Optional[List[str]] = None

class RecommendationResponse(BaseModel):
    """Recommendation response"""
    user_id: str
    recommendations: List[Dict]
    metadata: Dict

class RecommendationServer:
    """Recommendation serving with caching and business rules"""

    def __init__(self, redis_host: str, model_path: str):
        self.redis_client = redis.Redis(
            host=redis_host,
            decode_responses=False  # Binary for pickle
        )
        self.logger = logging.getLogger(__name__)

        # Load models (simplified - in production use model registry)
        with open(model_path, 'rb') as f:
            self.recommender = pickle.load(f)

        # Business rules
        self.min_inventory_threshold = 5
        self.min_margin_threshold = 0.15

    async def get_recommendations(
        self,
        request: RecommendationRequest
    ) -> RecommendationResponse:
        """Get personalized recommendations"""

        start_time = time.time()

        # Check cache first
        cache_key = f"recs:{request.user_id}:{request.num_recommendations}"
        cached = self.redis_client.get(cache_key)

        if cached:
            self.logger.info(f"Cache hit for user {request.user_id}")
            recommendations = pickle.loads(cached)
        else:
            # Generate recommendations
            recommendations = await self._generate_recommendations(request)

            # Cache for 5 minutes
            self.redis_client.setex(
                cache_key,
                300,
                pickle.dumps(recommendations)
            )

        # Apply business rules
        recommendations = self._apply_business_rules(recommendations)

        # Apply diversity
        recommendations = self._apply_diversity(recommendations)

        # Add metadata
        for rec in recommendations:
            rec['metadata'] = self._get_product_metadata(rec['product_id'])

        latency_ms = (time.time() - start_time) * 1000

        return RecommendationResponse(
            user_id=request.user_id,
            recommendations=recommendations[:request.num_recommendations],
            metadata={
                'latency_ms': latency_ms,
                'cache_hit': cached is not None,
                'timestamp': datetime.now().isoformat()
            }
        )

    async def _generate_recommendations(
        self,
        request: RecommendationRequest
    ) -> List[Dict]:
        """Generate recommendations using models"""

        # Get user features
        user_features = self._get_user_features(request.user_id)

        # Get raw recommendations
        raw_recs = self.recommender.recommend(
            user_id=request.user_id,
            user_features=user_features,
            n=50,  # Get more candidates for filtering
            exclude_products=request.exclude_products or []
        )

        # Convert to dict format
        recommendations = [
            {
                'product_id': product_id,
                'score': float(score),
                'reason': self._generate_explanation(product_id, user_features)
            }
            for product_id, score in raw_recs
        ]

        return recommendations

    def _apply_business_rules(self, recommendations: List[Dict]) -> List[Dict]:
        """Apply business constraints"""

        filtered = []

        for rec in recommendations:
            product_id = rec['product_id']

            # Check inventory
            inventory = self._get_inventory(product_id)
            if inventory < self.min_inventory_threshold:
                continue

            # Check margin
            margin = self._get_margin(product_id)
            if margin < self.min_margin_threshold:
                continue

            # Check if product is active
            if not self._is_product_active(product_id):
                continue

            filtered.append(rec)

        return filtered

    def _apply_diversity(self, recommendations: List[Dict]) -> List[Dict]:
        """Ensure category diversity"""

        if not recommendations:
            return recommendations

        diversified = [recommendations[0]]
        categories_seen = {self._get_category(recommendations[0]['product_id'])}

        for rec in recommendations[1:]:
            category = self._get_category(rec['product_id'])

            # Prefer products from unseen categories
            if category not in categories_seen or len(diversified) > 5:
                diversified.append(rec)
                categories_seen.add(category)

            if len(diversified) >= len(recommendations):
                break

        return diversified

    def _get_user_features(self, user_id: str) -> Dict:
        """Get user features from cache"""

        features_key = f"user:{user_id}:features"
        features = self.redis_client.get(features_key)

        if features:
            return pickle.loads(features)

        # Return default features for cold start
        return {
            'favorite_categories': [],
            'recent_views': [],
            'recent_purchases': []
        }

    def _get_product_metadata(self, product_id: str) -> Dict:
        """Get product metadata"""

        metadata_key = f"product:{product_id}:metadata"
        metadata = self.redis_client.get(metadata_key)

        if metadata:
            return pickle.loads(metadata)

        return {}

    def _get_inventory(self, product_id: str) -> int:
        """Get product inventory"""
        inventory = self.redis_client.get(f"product:{product_id}:inventory")
        return int(inventory) if inventory else 0

    def _get_margin(self, product_id: str) -> float:
        """Get product margin"""
        margin = self.redis_client.get(f"product:{product_id}:margin")
        return float(margin) if margin else 0.0

    def _is_product_active(self, product_id: str) -> bool:
        """Check if product is active"""
        active = self.redis_client.get(f"product:{product_id}:active")
        return active == b'1' if active else True

    def _get_category(self, product_id: str) -> str:
        """Get product category"""
        category = self.redis_client.get(f"product:{product_id}:category")
        return category.decode() if category else 'Unknown'

    def _generate_explanation(self, product_id: str, user_features: Dict) -> str:
        """Generate explanation for recommendation"""

        # Simple rule-based explanation
        if product_id in user_features.get('recent_views', []):
            return "Based on items you recently viewed"

        favorite_cat = user_features.get('favorite_categories', [None])[0]
        product_cat = self._get_category(product_id)

        if favorite_cat and product_cat == favorite_cat:
            return f"Popular in {product_cat}"

        return "Recommended for you"

# FastAPI endpoints
recommendation_server = None

@app.on_event("startup")
async def startup_event():
    global recommendation_server
    recommendation_server = RecommendationServer(
        redis_host='localhost',
        model_path='models/recommender.pkl'
    )

@app.post("/recommendations", response_model=RecommendationResponse)
async def get_recommendations(request: RecommendationRequest):
    """Get personalized recommendations"""
    try:
        return await recommendation_server.get_recommendations(request)
    except Exception as e:
        raise HTTPException(status_code=500, detail=str(e))

@app.get("/health")
async def health_check():
    """Health check endpoint"""
    return {"status": "healthy", "timestamp": datetime.now().isoformat()}

@app.get("/similar/{product_id}")
async def get_similar_products(
    product_id: str,
    n: int = Query(10, ge=1, le=50)
):
    """Get similar products"""
    try:
        similar = recommendation_server.recommender.cf_model.similar_items(
            product_id,
            n=n
        )

        return {
            "product_id": product_id,
            "similar_products": [
                {"product_id": pid, "score": float(score)}
                for pid, score in similar
            ]
        }
    except Exception as e:
        raise HTTPException(status_code=500, detail=str(e))

if __name__ == "__main__":
    import uvicorn
    uvicorn.run(app, host="0.0.0.0", port=8000)
\end{lstlisting}

\subsection{Results and Impact}

\textbf{Business Metrics (After 6 Months):}

\begin{itemize}
    \item \textbf{Click-Through Rate:} Increased from 2.3\% to 4.7\% (+104\%)
    \item \textbf{Conversion Rate:} Increased from 1.2\% to 2.1\% (+75\%)
    \item \textbf{Average Order Value:} Increased from \$47 to \$58 (+23\%)
    \item \textbf{Revenue Impact:} \$12M additional annual revenue
    \item \textbf{Customer Engagement:} Session duration increased by 35\%
\end{itemize}

\textbf{Technical Performance:}

\begin{itemize}
    \item \textbf{Latency:} p50: 45ms, p95: 78ms, p99: 92ms (target: <100ms)
    \item \textbf{Cache Hit Rate:} 87\% (significantly reduced compute costs)
    \item \textbf{Throughput:} 12,000 requests/second peak
    \item \textbf{Model Accuracy:} NDCG@10: 0.43, MAP@10: 0.38
\end{itemize}

\subsection{Lessons Learned}

\textbf{Key Insights:}

\begin{enumerate}
    \item \textbf{Hybrid Approach is Essential:} No single model performed well across all scenarios. Collaborative filtering excelled for popular users, content-based handled cold start, and deep learning captured complex patterns.

    \item \textbf{Real-Time Features Matter:} Adding real-time user session data improved relevance by 18\% compared to batch features alone.

    \item \textbf{Business Rules are Critical:} Pure ML recommendations often suggested out-of-stock or low-margin items. Business rule integration improved profitability by 31\%.

    \item \textbf{Diversity Prevents Boredom:} Users engaged 40\% longer when recommendations showed category diversity vs. showing only similar items.

    \item \textbf{Explanation Builds Trust:} Adding simple explanations increased click-through rates by 15\%.

    \item \textbf{Caching is Mandatory:} 87\% cache hit rate reduced serving costs by 6x while maintaining freshness.

    \item \textbf{A/B Testing Revealed Surprises:} Some seemingly good ML improvements (higher offline metrics) didn't translate to better business outcomes. Always test in production.

    \item \textbf{Cold Start is Hard:} 23\% of users were new each month. Content-based and trending items were essential for cold start.
\end{enumerate}

\textbf{Technical Challenges Overcome:}

\begin{itemize}
    \item \textbf{Real-time Feature Consistency:} Ensured training features matched serving features by using same feature store
    \item \textbf{Scale:} Handled 10x traffic during sales with autoscaling and precomputed recommendations
    \item \textbf{Model Updates:} Deployed new models daily without downtime using blue-green deployment
    \item \textbf{Data Quality:} Bot traffic and fraud required sophisticated filtering (12\% of events were invalid)
\end{itemize}

\section{Case Study 2: Financial Fraud Detection Pipeline}

\subsection{Business Context and Requirements}

\textbf{Company:} Digital payments platform processing \$50B annually\\
\textbf{Challenge:} Fraud losses of \$45M annually, manual review backlog\\
\textbf{Goal:} Reduce fraud losses while minimizing false positives

\textbf{Key Requirements:}

\begin{itemize}
    \item \textbf{Real-Time:} Decision in <100ms for transaction approval
    \item \textbf{High Recall:} Catch 95\%+ of fraudulent transactions
    \item \textbf{Low False Positives:} Keep legitimate user friction minimal (<1\% decline rate)
    \item \textbf{Explainability:} Provide reasons for fraud decisions (regulatory requirement)
    \item \textbf{Adaptability:} Detect new fraud patterns within hours
    \item \textbf{Security:} Protect sensitive financial data (PCI-DSS compliance)
\end{itemize}

\subsection{System Architecture}

\begin{figure}[h]
\centering
\begin{tikzpicture}[
    node distance=1.3cm,
    component/.style={rectangle, draw, fill=blue!20, text width=2.5cm, text centered, minimum height=0.9cm, font=\small},
    storage/.style={cylinder, draw, fill=green!20, text width=2.2cm, text centered, minimum height=0.9cm, aspect=0.3, font=\small},
    decision/.style={diamond, draw, fill=yellow!20, text width=2cm, text centered, minimum height=0.9cm, font=\small}
]

% Input
\node[component] (txn) {Transaction\\Request};

% Feature extraction
\node[component, below=of txn] (features) {Feature\\Extraction};
\node[storage, left=of features] (history) {User\\History};
\node[storage, right=of features] (network) {Network\\Graph};

% Models
\node[component, below=of features] (rules) {Rule\\Engine};
\node[component, left=of rules] (ml) {ML\\Models};
\node[component, right=of rules] (network_model) {Network\\Analysis};

% Decision
\node[decision, below=of rules] (decision) {Risk\\Score};

% Actions
\node[component, below left=of decision] (approve) {Approve};
\node[component, below=of decision] (review) {Manual\\Review};
\node[component, below right=of decision] (block) {Block};

% Feedback
\node[storage, below=of review] (feedback) {Feedback\\Loop};

% Arrows
\draw[->] (txn) -- (features);
\draw[->] (history) -- (features);
\draw[->] (network) -- (features);
\draw[->] (features) -- (rules);
\draw[->] (features) -- (ml);
\draw[->] (features) -- (network_model);
\draw[->] (rules) -- (decision);
\draw[->] (ml) -- (decision);
\draw[->] (network_model) -- (decision);
\draw[->] (decision) -- node[left] {Low} (approve);
\draw[->] (decision) -- node[above] {Medium} (review);
\draw[->] (decision) -- node[right] {High} (block);
\draw[->] (review) -- (feedback);
\draw[->] (feedback) -| (ml);

\end{tikzpicture}
\caption{Fraud Detection System Architecture}
\end{figure}

\subsection{Feature Engineering and Real-Time Processing}

\begin{lstlisting}[language=Python, caption=Real-Time Fraud Detection Features]
#!/usr/bin/env python3
"""
Financial Fraud Detection - Real-Time Feature Engineering
Extract transaction features in <50ms for real-time fraud detection
"""

from dataclasses import dataclass
from typing import Dict, List, Optional
from datetime import datetime, timedelta
import redis
import numpy as np
from geopy.distance import geodesic
import hashlib

@dataclass
class Transaction:
    """Transaction data"""
    transaction_id: str
    user_id: str
    merchant_id: str
    amount: float
    currency: str
    timestamp: datetime
    location: tuple  # (latitude, longitude)
    device_id: str
    ip_address: str
    payment_method: str
    metadata: Dict

@dataclass
class FraudFeatures:
    """Computed fraud detection features"""
    # User velocity features
    txns_last_hour: int
    txns_last_24h: int
    amount_last_hour: float
    amount_last_24h: float
    unique_merchants_last_hour: int

    # Geographic features
    distance_from_last_txn_km: float
    velocity_km_per_hour: float
    new_country: bool
    high_risk_country: bool

    # Device and identity
    new_device: bool
    device_change_frequency: float
    ip_reputation_score: float

    # Merchant features
    merchant_risk_score: float
    first_time_merchant: bool

    # Amount features
    amount_deviation_from_avg: float
    round_amount: bool
    amount_percentile: float

    # Time features
    hour_of_day: int
    day_of_week: int
    unusual_hour: bool

    # Network features
    shared_device_users: int
    shared_ip_users: int

class FraudFeatureExtractor:
    """Extract fraud detection features in real-time"""

    def __init__(self, redis_host: str):
        self.redis_client = redis.Redis(host=redis_host, decode_responses=True)

        # High-risk countries (example list)
        self.high_risk_countries = {'XX', 'YY', 'ZZ'}

        # Typical transaction hours by user
        self.typical_hours_start = 6
        self.typical_hours_end = 23

    def extract_features(self, transaction: Transaction) -> FraudFeatures:
        """Extract all features for fraud detection"""

        # User velocity features
        velocity_features = self._compute_velocity_features(transaction)

        # Geographic features
        geo_features = self._compute_geographic_features(transaction)

        # Device and identity features
        device_features = self._compute_device_features(transaction)

        # Merchant features
        merchant_features = self._compute_merchant_features(transaction)

        # Amount features
        amount_features = self._compute_amount_features(transaction)

        # Time features
        time_features = self._compute_time_features(transaction)

        # Network features
        network_features = self._compute_network_features(transaction)

        # Combine all features
        features = FraudFeatures(
            **velocity_features,
            **geo_features,
            **device_features,
            **merchant_features,
            **amount_features,
            **time_features,
            **network_features
        )

        # Update historical data for future transactions
        self._update_user_history(transaction)

        return features

    def _compute_velocity_features(self, txn: Transaction) -> Dict:
        """Compute transaction velocity features"""

        user_key = f"user:{txn.user_id}"

        # Count transactions in time windows
        now = txn.timestamp.timestamp()
        hour_ago = now - 3600
        day_ago = now - 86400

        # Get recent transactions
        recent_txns = self.redis_client.zrangebyscore(
            f"{user_key}:txns",
            hour_ago,
            now
        )

        txns_last_hour = len(recent_txns)

        recent_txns_24h = self.redis_client.zrangebyscore(
            f"{user_key}:txns",
            day_ago,
            now
        )

        txns_last_24h = len(recent_txns_24h)

        # Sum amounts
        amount_last_hour = 0.0
        for txn_data in recent_txns:
            amount = float(txn_data.split(':')[1])
            amount_last_hour += amount

        amount_last_24h = 0.0
        for txn_data in recent_txns_24h:
            amount = float(txn_data.split(':')[1])
            amount_last_24h += amount

        # Count unique merchants
        merchants_last_hour = set()
        for txn_data in recent_txns:
            merchant = txn_data.split(':')[2]
            merchants_last_hour.add(merchant)

        return {
            'txns_last_hour': txns_last_hour,
            'txns_last_24h': txns_last_24h,
            'amount_last_hour': amount_last_hour,
            'amount_last_24h': amount_last_24h,
            'unique_merchants_last_hour': len(merchants_last_hour)
        }

    def _compute_geographic_features(self, txn: Transaction) -> Dict:
        """Compute geographic fraud indicators"""

        user_key = f"user:{txn.user_id}"

        # Get last transaction location
        last_location_data = self.redis_client.get(f"{user_key}:last_location")

        if last_location_data:
            parts = last_location_data.split(',')
            last_lat, last_lon, last_time_str = float(parts[0]), float(parts[1]), parts[2]
            last_time = datetime.fromisoformat(last_time_str)

            # Calculate distance
            distance_km = geodesic(
                (last_lat, last_lon),
                txn.location
            ).kilometers

            # Calculate velocity (km/h)
            time_diff_hours = (txn.timestamp - last_time).total_seconds() / 3600
            velocity = distance_km / time_diff_hours if time_diff_hours > 0 else 0

            # Flag impossible travel (> 1000 km/h)
            impossible_travel = velocity > 1000

        else:
            distance_km = 0
            velocity = 0
            impossible_travel = False

        # Check if new country
        user_countries = self.redis_client.smembers(f"{user_key}:countries")
        country = self._get_country_from_location(txn.location)
        new_country = country not in user_countries

        high_risk_country = country in self.high_risk_countries

        return {
            'distance_from_last_txn_km': distance_km,
            'velocity_km_per_hour': velocity,
            'new_country': new_country,
            'high_risk_country': high_risk_country
        }

    def _compute_device_features(self, txn: Transaction) -> Dict:
        """Compute device and identity features"""

        user_key = f"user:{txn.user_id}"

        # Check if new device
        user_devices = self.redis_client.smembers(f"{user_key}:devices")
        new_device = txn.device_id not in user_devices

        # Calculate device change frequency
        device_changes = int(self.redis_client.get(f"{user_key}:device_changes") or 0)
        total_txns = int(self.redis_client.get(f"{user_key}:total_txns") or 1)
        device_change_frequency = device_changes / total_txns

        # IP reputation (simplified - in production use threat intel services)
        ip_reputation_score = self._get_ip_reputation(txn.ip_address)

        return {
            'new_device': new_device,
            'device_change_frequency': device_change_frequency,
            'ip_reputation_score': ip_reputation_score
        }

    def _compute_merchant_features(self, txn: Transaction) -> Dict:
        """Compute merchant risk features"""

        user_key = f"user:{txn.user_id}"
        merchant_key = f"merchant:{txn.merchant_id}"

        # Check if first time with merchant
        user_merchants = self.redis_client.smembers(f"{user_key}:merchants")
        first_time_merchant = txn.merchant_id not in user_merchants

        # Get merchant risk score
        merchant_risk = float(self.redis_client.get(f"{merchant_key}:risk_score") or 0.5)

        return {
            'merchant_risk_score': merchant_risk,
            'first_time_merchant': first_time_merchant
        }

    def _compute_amount_features(self, txn: Transaction) -> Dict:
        """Compute amount-based features"""

        user_key = f"user:{txn.user_id}"

        # Get user's historical average
        avg_amount = float(self.redis_client.get(f"{user_key}:avg_amount") or 100.0)

        # Deviation from average
        deviation = abs(txn.amount - avg_amount) / avg_amount if avg_amount > 0 else 0

        # Check if round amount (potential test transaction)
        round_amount = txn.amount % 10 == 0 and txn.amount >= 100

        # Get percentile of this amount
        # Simplified - in production, maintain histogram
        amount_percentile = min(txn.amount / (avg_amount * 3), 1.0)

        return {
            'amount_deviation_from_avg': deviation,
            'round_amount': round_amount,
            'amount_percentile': amount_percentile
        }

    def _compute_time_features(self, txn: Transaction) -> Dict:
        """Compute time-based features"""

        hour = txn.timestamp.hour
        day = txn.timestamp.weekday()

        # Check if unusual hour
        unusual_hour = hour < self.typical_hours_start or hour > self.typical_hours_end

        return {
            'hour_of_day': hour,
            'day_of_week': day,
            'unusual_hour': unusual_hour
        }

    def _compute_network_features(self, txn: Transaction) -> Dict:
        """Compute network-based features (shared devices/IPs)"""

        # How many users share this device?
        device_users = self.redis_client.scard(f"device:{txn.device_id}:users")

        # How many users share this IP?
        ip_users = self.redis_client.scard(f"ip:{txn.ip_address}:users")

        return {
            'shared_device_users': device_users,
            'shared_ip_users': ip_users
        }

    def _update_user_history(self, txn: Transaction):
        """Update user history after transaction"""

        user_key = f"user:{txn.user_id}"

        # Add transaction to sorted set (by timestamp)
        txn_data = f"{txn.transaction_id}:{txn.amount}:{txn.merchant_id}"
        self.redis_client.zadd(
            f"{user_key}:txns",
            {txn_data: txn.timestamp.timestamp()}
        )

        # Keep only last 30 days
        cutoff = (txn.timestamp - timedelta(days=30)).timestamp()
        self.redis_client.zremrangebyscore(f"{user_key}:txns", 0, cutoff)

        # Update location
        location_data = f"{txn.location[0]},{txn.location[1]},{txn.timestamp.isoformat()}"
        self.redis_client.set(f"{user_key}:last_location", location_data)

        # Update country set
        country = self._get_country_from_location(txn.location)
        self.redis_client.sadd(f"{user_key}:countries", country)

        # Update device set
        if txn.device_id not in self.redis_client.smembers(f"{user_key}:devices"):
            self.redis_client.incr(f"{user_key}:device_changes")
        self.redis_client.sadd(f"{user_key}:devices", txn.device_id)

        # Update merchant set
        self.redis_client.sadd(f"{user_key}:merchants", txn.merchant_id)

        # Update total transactions
        self.redis_client.incr(f"{user_key}:total_txns")

        # Update average amount (moving average)
        current_avg = float(self.redis_client.get(f"{user_key}:avg_amount") or txn.amount)
        count = int(self.redis_client.get(f"{user_key}:total_txns") or 1)
        new_avg = (current_avg * (count - 1) + txn.amount) / count
        self.redis_client.set(f"{user_key}:avg_amount", new_avg)

        # Network features
        self.redis_client.sadd(f"device:{txn.device_id}:users", txn.user_id)
        self.redis_client.sadd(f"ip:{txn.ip_address}:users", txn.user_id)

    def _get_country_from_location(self, location: tuple) -> str:
        """Get country code from lat/lon"""
        # Simplified - in production use reverse geocoding service
        return "US"

    def _get_ip_reputation(self, ip_address: str) -> float:
        """Get IP reputation score"""
        # Simplified - in production integrate with threat intelligence
        return 0.8

# Example usage
if __name__ == "__main__":
    extractor = FraudFeatureExtractor(redis_host='localhost')

    transaction = Transaction(
        transaction_id='txn_12345',
        user_id='user_789',
        merchant_id='merchant_456',
        amount=1250.00,
        currency='USD',
        timestamp=datetime.now(),
        location=(40.7128, -74.0060),  # NYC
        device_id='device_abc123',
        ip_address='192.168.1.1',
        payment_method='credit_card',
        metadata={}
    )

    features = extractor.extract_features(transaction)

    print(f"Transactions last hour: {features.txns_last_hour}")
    print(f"Amount deviation: {features.amount_deviation_from_avg:.2f}")
    print(f"New device: {features.new_device}")
    print(f"Travel velocity: {features.velocity_km_per_hour:.1f} km/h")
    print(f"Merchant risk: {features.merchant_risk_score:.2f}")
\end{lstlisting}

\subsection{Fraud Detection Models}

\begin{lstlisting}[language=Python, caption=Ensemble Fraud Detection Models]
#!/usr/bin/env python3
"""
Fraud Detection Models
Ensemble of models for robust fraud detection
"""

import xgboost as xgb
import lightgbm as lgb
from sklearn.ensemble import IsolationForest
from sklearn.neural_network import MLPClassifier
import numpy as np
import pandas as pd
from typing import List, Dict, Tuple
import shap

class RuleEngine:
    """Hard rules for known fraud patterns"""

    def __init__(self):
        self.rules = []
        self._define_rules()

    def _define_rules(self):
        """Define fraud detection rules"""

        # Rule 1: Impossible travel
        self.rules.append({
            'name': 'impossible_travel',
            'condition': lambda f: f.velocity_km_per_hour > 1000,
            'risk_score': 0.95,
            'reason': 'Impossible travel speed detected'
        })

        # Rule 2: High velocity in short time
        self.rules.append({
            'name': 'high_velocity',
            'condition': lambda f: f.txns_last_hour > 10,
            'risk_score': 0.85,
            'reason': 'Unusually high transaction velocity'
        })

        # Rule 3: Large amount on new device
        self.rules.append({
            'name': 'new_device_large_amount',
            'condition': lambda f: f.new_device and f.amount_percentile > 0.95,
            'risk_score': 0.75,
            'reason': 'Large transaction from new device'
        })

        # Rule 4: High-risk merchant and country
        self.rules.append({
            'name': 'high_risk_combination',
            'condition': lambda f: f.high_risk_country and f.merchant_risk_score > 0.7,
            'risk_score': 0.80,
            'reason': 'High-risk country and merchant combination'
        })

        # Rule 5: Round amount test transaction
        self.rules.append({
            'name': 'test_transaction',
            'condition': lambda f: f.round_amount and f.first_time_merchant,
            'risk_score': 0.65,
            'reason': 'Potential test transaction pattern'
        })

    def evaluate(self, features: FraudFeatures) -> Tuple[float, List[str]]:
        """Evaluate rules and return risk score"""

        max_risk = 0.0
        triggered_rules = []

        for rule in self.rules:
            if rule['condition'](features):
                max_risk = max(max_risk, rule['risk_score'])
                triggered_rules.append(rule['reason'])

        return max_risk, triggered_rules

class FraudDetectionModel:
    """ML model for fraud detection"""

    def __init__(self):
        # Primary model: XGBoost
        self.xgb_model = xgb.XGBClassifier(
            max_depth=6,
            learning_rate=0.1,
            n_estimators=200,
            objective='binary:logistic',
            scale_pos_weight=50,  # Handle class imbalance
            random_state=42
        )

        # Secondary model: LightGBM
        self.lgb_model = lgb.LGBMClassifier(
            num_leaves=31,
            learning_rate=0.05,
            n_estimators=200,
            class_weight='balanced',
            random_state=42
        )

        # Anomaly detection: Isolation Forest
        self.anomaly_model = IsolationForest(
            contamination=0.01,
            random_state=42
        )

        # Feature names
        self.feature_names = None

        # SHAP explainer
        self.explainer = None

    def train(self, X: pd.DataFrame, y: np.ndarray):
        """Train all models"""

        self.feature_names = X.columns.tolist()

        print("Training XGBoost...")
        self.xgb_model.fit(X, y)

        print("Training LightGBM...")
        self.lgb_model.fit(X, y)

        print("Training Isolation Forest...")
        self.anomaly_model.fit(X)

        # Initialize SHAP explainer
        print("Creating SHAP explainer...")
        self.explainer = shap.TreeExplainer(self.xgb_model)

        print("Training complete")

    def predict_proba(self, X: pd.DataFrame) -> np.ndarray:
        """Predict fraud probability (ensemble)"""

        # XGBoost predictions
        xgb_proba = self.xgb_model.predict_proba(X)[:, 1]

        # LightGBM predictions
        lgb_proba = self.lgb_model.predict_proba(X)[:, 1]

        # Isolation Forest anomaly scores (convert to probability)
        anomaly_scores = self.anomaly_model.score_samples(X)
        anomaly_proba = 1 / (1 + np.exp(anomaly_scores))  # Sigmoid transformation

        # Ensemble: weighted average
        ensemble_proba = (
            0.5 * xgb_proba +
            0.4 * lgb_proba +
            0.1 * anomaly_proba
        )

        return ensemble_proba

    def explain(self, X: pd.DataFrame) -> Dict:
        """Generate SHAP explanations"""

        shap_values = self.explainer.shap_values(X)

        # Get top contributing features
        feature_importance = {}
        for i, feature_name in enumerate(self.feature_names):
            feature_importance[feature_name] = float(np.abs(shap_values[0][i]))

        # Sort by importance
        sorted_features = sorted(
            feature_importance.items(),
            key=lambda x: x[1],
            reverse=True
        )

        return {
            'top_features': sorted_features[:5],
            'shap_values': shap_values[0].tolist()
        }

class FraudDetectionEnsemble:
    """Complete fraud detection system"""

    def __init__(self):
        self.rule_engine = RuleEngine()
        self.ml_model = FraudDetectionModel()

        # Thresholds
        self.block_threshold = 0.90
        self.review_threshold = 0.50

    def predict(
        self,
        features: FraudFeatures
    ) -> Dict:
        """Predict fraud risk"""

        # Convert features to DataFrame
        feature_dict = vars(features)
        X = pd.DataFrame([feature_dict])

        # Rule-based score
        rule_score, triggered_rules = self.rule_engine.evaluate(features)

        # ML model score
        ml_score = self.ml_model.predict_proba(X)[0]

        # Combined score (max of rule and ML)
        final_score = max(rule_score, ml_score * 0.9)  # Slight discount on ML

        # Decision
        if final_score >= self.block_threshold:
            decision = 'BLOCK'
        elif final_score >= self.review_threshold:
            decision = 'REVIEW'
        else:
            decision = 'APPROVE'

        # Generate explanation
        if triggered_rules:
            explanation = triggered_rules
        else:
            # Use SHAP for ML explanation
            shap_explanation = self.ml_model.explain(X)
            explanation = [
                f"{feature}: {importance:.3f}"
                for feature, importance in shap_explanation['top_features']
            ]

        return {
            'decision': decision,
            'risk_score': float(final_score),
            'rule_score': float(rule_score),
            'ml_score': float(ml_score),
            'explanation': explanation,
            'triggered_rules': triggered_rules
        }

    def update_with_feedback(
        self,
        transaction_id: str,
        features: pd.DataFrame,
        label: int
    ):
        """Update model with feedback (online learning)"""

        # In production, this would trigger retraining pipeline
        # For now, just log for batch retraining

        print(f"Feedback received for {transaction_id}: {'Fraud' if label == 1 else 'Legitimate'}")

# Example usage
if __name__ == "__main__":
    # Load training data
    train_data = pd.read_csv('fraud_training_data.csv')
    X_train = train_data.drop('is_fraud', axis=1)
    y_train = train_data['is_fraud']

    # Train ensemble
    ensemble = FraudDetectionEnsemble()
    ensemble.ml_model.train(X_train, y_train)

    # Predict on new transaction
    features = FraudFeatures(
        txns_last_hour=15,
        txns_last_24h=45,
        amount_last_hour=5000,
        amount_last_24h=12000,
        unique_merchants_last_hour=8,
        distance_from_last_txn_km=2500,
        velocity_km_per_hour=1200,  # Suspicious
        new_country=True,
        high_risk_country=False,
        new_device=True,
        device_change_frequency=0.3,
        ip_reputation_score=0.6,
        merchant_risk_score=0.4,
        first_time_merchant=True,
        amount_deviation_from_avg=4.5,
        round_amount=True,
        amount_percentile=0.95,
        hour_of_day=3,
        day_of_week=2,
        unusual_hour=True,
        shared_device_users=1,
        shared_ip_users=2
    )

    prediction = ensemble.predict(features)

    print(f"Decision: {prediction['decision']}")
    print(f"Risk Score: {prediction['risk_score']:.3f}")
    print(f"Explanation: {prediction['explanation']}")
\end{lstlisting}

\subsection{Results and Impact}

\textbf{Fraud Prevention Metrics:}

\begin{itemize}
    \item \textbf{Fraud Detection Rate:} 96.5\% (up from 87\%)
    \item \textbf{False Positive Rate:} 0.8\% (down from 2.3\%)
    \item \textbf{Fraud Losses:} Reduced from \$45M to \$8M annually (82\% reduction)
    \item \textbf{Manual Review Queue:} Reduced by 60\%
    \item \textbf{Genuine User Friction:} Decreased by 65\%
\end{itemize}

\textbf{Operational Efficiency:}

\begin{itemize}
    \item \textbf{Decision Latency:} p50: 42ms, p99: 87ms
    \item \textbf{Throughput:} 50,000 transactions/second
    \item \textbf{Model Update Frequency:} Retrained daily, deployed hourly
    \item \textbf{Investigation Time:} Reduced from 15 minutes to 3 minutes per case
\end{itemize}

\subsection{Lessons Learned}

\textbf{Key Insights:}

\begin{enumerate}
    \item \textbf{Rules + ML Better Than Either Alone:} Known fraud patterns caught by rules, novel patterns by ML. Hybrid approach reduced fraud by 82\% vs. 61\% for ML-only.

    \item \textbf{Real-Time Features Critical:} Transaction velocity and impossible travel features had highest predictive power. Batch features alone missed 23\% of fraud.

    \item \textbf{Explainability is Non-Negotiable:} Regulatory requirements demanded explanations. SHAP values satisfied auditors and helped fraud analysts.

    \item \textbf{Class Imbalance is Severe:} Only 0.1\% of transactions were fraud. Required aggressive class weighting and careful metric selection (precision-recall over accuracy).

    \item \textbf{Fraud Evolves Rapidly:} New patterns emerged within days. Daily retraining and hourly deployment were essential.

    \item \textbf{False Positives Cost Real Money:} Each false decline cost \$50 in lost revenue and customer friction. Optimizing precision was as important as recall.

    \item \textbf{Network Features Surprised Us:} Shared device/IP features caught 15\% of fraud that other features missed. Fraudsters often share infrastructure.

    \item \textbf{A/B Testing Required Care:} Couldn't A/B test fraud models directly (ethical issues). Used historical replay and shadow mode instead.
\end{enumerate}

\section{Case Study 3: Healthcare Diagnostic System}

\subsection{Business Context and Requirements}

\textbf{Organization:} Healthcare network with 50 hospitals\\
\textbf{Challenge:} Radiologist shortage, diagnostic delays, missed findings\\
\textbf{Goal:} AI-assisted diagnosis to improve accuracy and speed

\textbf{Key Requirements:}

\begin{itemize}
    \item \textbf{Safety:} FDA clearance, >95\% sensitivity for critical findings
    \item \textbf{Interpretability:} Visual explanations for radiologists
    \item \textbf{Integration:} Seamless PACS/EMR integration
    \item \textbf{Performance:} Process images in <30 seconds
    \item \textbf{Privacy:} HIPAA compliance, de-identification
    \item \textbf{Reliability:} 99.9\% uptime for clinical operations
\end{itemize}

\subsection{System Architecture}

\begin{lstlisting}[language=Python, caption=Medical Image Analysis Pipeline]
#!/usr/bin/env python3
"""
Healthcare Diagnostic System - Medical Image Analysis
AI-assisted diagnosis for chest X-rays with FDA compliance
"""

import torch
import torch.nn as nn
import torchvision.transforms as transforms
from torchvision.models import densenet121
import cv2
import numpy as np
from typing import List, Dict, Tuple, Optional
from dataclasses import dataclass
import pydicom
from PIL import Image
import gradcam
from datetime import datetime

@dataclass
class Finding:
    """Medical finding"""
    condition: str
    probability: float
    confidence_interval: Tuple[float, float]
    severity: str  # normal, mild, moderate, severe
    location: Optional[Dict] = None  # Bounding box if available

@dataclass
class DiagnosticReport:
    """Complete diagnostic report"""
    study_id: str
    patient_id_hash: str
    image_type: str
    findings: List[Finding]
    overall_assessment: str
    urgent_findings: List[str]
    attention_map: np.ndarray
    processing_time_ms: float
    model_version: str
    timestamp: datetime

class MedicalImagePreprocessor:
    """Preprocess medical images for analysis"""

    def __init__(self):
        self.target_size = (224, 224)

        self.transform = transforms.Compose([
            transforms.Resize(self.target_size),
            transforms.ToTensor(),
            transforms.Normalize(
                mean=[0.485, 0.456, 0.406],
                std=[0.229, 0.224, 0.225]
            )
        ])

    def load_dicom(self, dicom_path: str) -> Tuple[np.ndarray, Dict]:
        """Load and parse DICOM file"""

        # Read DICOM
        dicom = pydicom.dcmread(dicom_path)

        # Extract pixel data
        image = dicom.pixel_array

        # Apply DICOM windowing
        image = self._apply_window(image, dicom)

        # Extract metadata
        metadata = {
            'study_id': str(dicom.StudyInstanceUID),
            'series_id': str(dicom.SeriesInstanceUID),
            'patient_age': dicom.get('PatientAge', 'Unknown'),
            'patient_sex': dicom.get('PatientSex', 'Unknown'),
            'study_date': str(dicom.StudyDate),
            'modality': dicom.Modality,
            'view_position': dicom.get('ViewPosition', 'Unknown')
        }

        return image, metadata

    def _apply_window(self, image: np.ndarray, dicom) -> np.ndarray:
        """Apply DICOM windowing"""

        window_center = dicom.get('WindowCenter', 40)
        window_width = dicom.get('WindowWidth', 400)

        if isinstance(window_center, pydicom.multival.MultiValue):
            window_center = window_center[0]
        if isinstance(window_width, pydicom.multival.MultiValue):
            window_width = window_width[0]

        min_value = window_center - window_width / 2
        max_value = window_center + window_width / 2

        # Apply window
        image = np.clip(image, min_value, max_value)

        # Normalize to 0-255
        image = ((image - min_value) / (max_value - min_value) * 255).astype(np.uint8)

        return image

    def preprocess(self, image: np.ndarray) -> torch.Tensor:
        """Preprocess image for model"""

        # Convert to 3-channel if grayscale
        if len(image.shape) == 2:
            image = cv2.cvtColor(image, cv2.COLOR_GRAY2RGB)

        # Convert to PIL Image
        image = Image.fromarray(image)

        # Apply transforms
        tensor = self.transform(image)

        return tensor.unsqueeze(0)  # Add batch dimension

class ChestXRayModel(nn.Module):
    """Multi-label classification model for chest X-rays"""

    def __init__(self, num_classes: int = 14):
        super().__init__()

        # Base model: DenseNet121 pretrained on ImageNet
        self.densenet = densenet121(pretrained=True)

        # Replace classifier
        num_features = self.densenet.classifier.in_features
        self.densenet.classifier = nn.Sequential(
            nn.Linear(num_features, 512),
            nn.ReLU(),
            nn.Dropout(0.3),
            nn.Linear(512, num_classes)
        )

        # Condition names (ChestX-ray14 dataset)
        self.conditions = [
            'Atelectasis',
            'Cardiomegaly',
            'Effusion',
            'Infiltration',
            'Mass',
            'Nodule',
            'Pneumonia',
            'Pneumothorax',
            'Consolidation',
            'Edema',
            'Emphysema',
            'Fibrosis',
            'Pleural_Thickening',
            'Hernia'
        ]

        # Critical findings requiring urgent attention
        self.critical_findings = {
            'Pneumothorax',
            'Mass',
            'Pneumonia'
        }

    def forward(self, x):
        """Forward pass"""
        return self.densenet(x)

    def predict_with_uncertainty(
        self,
        x: torch.Tensor,
        num_samples: int = 20
    ) -> Tuple[torch.Tensor, torch.Tensor]:
        """Predict with uncertainty estimation using MC Dropout"""

        self.train()  # Enable dropout

        predictions = []
        for _ in range(num_samples):
            with torch.no_grad():
                logits = self.forward(x)
                probs = torch.sigmoid(logits)
                predictions.append(probs)

        predictions = torch.stack(predictions)

        # Mean and standard deviation
        mean_pred = predictions.mean(dim=0)
        std_pred = predictions.std(dim=0)

        self.eval()

        return mean_pred, std_pred

class DiagnosticSystem:
    """Complete diagnostic system"""

    def __init__(self, model_path: str, device: str = 'cuda'):
        self.device = device
        self.preprocessor = MedicalImagePreprocessor()

        # Load model
        self.model = ChestXRayModel()
        self.model.load_state_dict(torch.load(model_path))
        self.model.to(device)
        self.model.eval()

        # GradCAM for visualization
        self.gradcam = gradcam.GradCAM(self.model, self.model.densenet.features)

        # Calibration (for better probability estimates)
        self.calibration_parameters = self._load_calibration()

    def analyze_image(
        self,
        dicom_path: str,
        generate_attention_map: bool = True
    ) -> DiagnosticReport:
        """Analyze medical image and generate report"""

        start_time = datetime.now()

        # Load and preprocess image
        image, metadata = self.preprocessor.load_dicom(dicom_path)
        tensor = self.preprocessor.preprocess(image).to(self.device)

        # Predict with uncertainty
        probabilities, uncertainties = self.model.predict_with_uncertainty(tensor)
        probabilities = probabilities.cpu().numpy()[0]
        uncertainties = uncertainties.cpu().numpy()[0]

        # Apply calibration
        probabilities = self._calibrate_probabilities(probabilities)

        # Generate findings
        findings = []
        urgent_findings = []

        for i, condition in enumerate(self.model.conditions):
            prob = float(probabilities[i])
            unc = float(uncertainties[i])

            # Only report significant findings (>10% probability)
            if prob > 0.10:
                # Confidence interval (95%)
                ci_lower = max(0, prob - 1.96 * unc)
                ci_upper = min(1, prob + 1.96 * unc)

                # Determine severity
                if prob < 0.3:
                    severity = 'mild'
                elif prob < 0.6:
                    severity = 'moderate'
                else:
                    severity = 'severe'

                finding = Finding(
                    condition=condition,
                    probability=prob,
                    confidence_interval=(ci_lower, ci_upper),
                    severity=severity
                )

                findings.append(finding)

                # Check if urgent
                if condition in self.model.critical_findings and prob > 0.5:
                    urgent_findings.append(condition)

        # Sort findings by probability
        findings.sort(key=lambda f: f.probability, reverse=True)

        # Generate attention map
        attention_map = None
        if generate_attention_map and findings:
            attention_map = self._generate_attention_map(tensor, image)

        # Overall assessment
        if urgent_findings:
            assessment = f"URGENT: Possible {', '.join(urgent_findings)} detected"
        elif findings:
            assessment = f"Findings detected: {', '.join([f.condition for f in findings[:3]])}"
        else:
            assessment = "No significant findings"

        processing_time = (datetime.now() - start_time).total_seconds() * 1000

        return DiagnosticReport(
            study_id=metadata['study_id'],
            patient_id_hash=self._hash_patient_id(metadata.get('patient_id', '')),
            image_type=f"{metadata['modality']} {metadata['view_position']}",
            findings=findings,
            overall_assessment=assessment,
            urgent_findings=urgent_findings,
            attention_map=attention_map,
            processing_time_ms=processing_time,
            model_version='v2.1.0',
            timestamp=datetime.now()
        )

    def _generate_attention_map(
        self,
        tensor: torch.Tensor,
        original_image: np.ndarray
    ) -> np.ndarray:
        """Generate GradCAM attention map"""

        # Generate GradCAM
        cam = self.gradcam.generate(tensor)

        # Resize to original image size
        cam = cv2.resize(cam, (original_image.shape[1], original_image.shape[0]))

        # Create heatmap overlay
        heatmap = cv2.applyColorMap(np.uint8(255 * cam), cv2.COLORMAP_JET)

        # Convert original image to 3-channel
        if len(original_image.shape) == 2:
            original_image = cv2.cvtColor(original_image, cv2.COLOR_GRAY2RGB)

        # Overlay
        overlay = cv2.addWeighted(original_image, 0.6, heatmap, 0.4, 0)

        return overlay

    def _calibrate_probabilities(self, probabilities: np.ndarray) -> np.ndarray:
        """Apply temperature scaling for calibration"""

        # Temperature parameter (learned from validation set)
        temperature = self.calibration_parameters.get('temperature', 1.5)

        # Apply temperature scaling
        logits = np.log(probabilities / (1 - probabilities + 1e-10))
        calibrated_logits = logits / temperature
        calibrated_probs = 1 / (1 + np.exp(-calibrated_logits))

        return calibrated_probs

    def _load_calibration(self) -> Dict:
        """Load calibration parameters"""
        # In production, load from file
        return {'temperature': 1.5}

    def _hash_patient_id(self, patient_id: str) -> str:
        """Hash patient ID for privacy"""
        import hashlib
        return hashlib.sha256(patient_id.encode()).hexdigest()[:16]

    def generate_clinical_report(self, report: DiagnosticReport) -> str:
        """Generate clinical report text"""

        lines = []
        lines.append("=" * 60)
        lines.append("AI-ASSISTED DIAGNOSTIC REPORT")
        lines.append("=" * 60)
        lines.append(f"Study ID: {report.study_id}")
        lines.append(f"Image Type: {report.image_type}")
        lines.append(f"Analysis Date: {report.timestamp.strftime('%Y-%m-%d %H:%M:%S')}")
        lines.append(f"Model Version: {report.model_version}")
        lines.append("")

        if report.urgent_findings:
            lines.append("*** URGENT FINDINGS ***")
            for finding in report.urgent_findings:
                lines.append(f"  - {finding}")
            lines.append("")

        lines.append("FINDINGS:")
        if report.findings:
            for finding in report.findings:
                lines.append(f"\n  {finding.condition}:")
                lines.append(f"    Probability: {finding.probability:.1%}")
                lines.append(f"    Confidence Interval: [{finding.confidence_interval[0]:.1%}, {finding.confidence_interval[1]:.1%}]")
                lines.append(f"    Severity: {finding.severity.upper()}")
        else:
            lines.append("  No significant findings detected")

        lines.append(f"\nOVERALL ASSESSMENT:")
        lines.append(f"  {report.overall_assessment}")

        lines.append(f"\nProcessing Time: {report.processing_time_ms:.1f}ms")

        lines.append("\n" + "=" * 60)
        lines.append("DISCLAIMER: This is an AI-assisted analysis tool.")
        lines.append("Final diagnosis must be confirmed by a qualified radiologist.")
        lines.append("=" * 60)

        return "\n".join(lines)

# Example usage
if __name__ == "__main__":
    system = DiagnosticSystem(
        model_path='models/chestxray_densenet121.pth',
        device='cuda' if torch.cuda.is_available() else 'cpu'
    )

    # Analyze chest X-ray
    report = system.analyze_image('patient_xray.dcm')

    # Print clinical report
    print(system.generate_clinical_report(report))

    # Save attention map
    if report.attention_map is not None:
        cv2.imwrite('attention_map.png', report.attention_map)
\end{lstlisting}

\subsection{Results and Impact}

\textbf{Clinical Outcomes:}

\begin{itemize}
    \item \textbf{Diagnostic Accuracy:} 94.2\% sensitivity, 89.7\% specificity
    \item \textbf{Radiologist Efficiency:} Reading time reduced from 6 min to 3.5 min per study
    \item \textbf{Critical Finding Detection:} Pneumothorax detection improved from 92\% to 97\%
    \item \textbf{Missed Findings:} Reduced by 31\% through AI second-read
    \item \textbf{Throughput:} Enabled 40\% increase in studies read per day
\end{itemize}

\textbf{Operational Benefits:}

\begin{itemize}
    \item \textbf{Wait Time:} Reduced from 48 hours to 12 hours average
    \item \textbf{After-Hours Coverage:} AI triage enabled remote radiologist efficiency
    \item \textbf{Teaching Tool:} Used for resident training with attention maps
\end{itemize}

\subsection{Lessons Learned}

\begin{enumerate}
    \item \textbf{FDA Clearance Required Partnership:} Clinical validation and regulatory submission took 18 months. Started early with FDA pre-submission meeting.

    \item \textbf{Calibration is Critical:} Raw model probabilities were overconfident. Temperature scaling improved reliability and radiologist trust.

    \item \textbf{Attention Maps Changed Everything:} GradCAM visualizations built trust with radiologists. Without interpretability, adoption would have failed.

    \item \textbf{Integration Matters Most:} PACS integration was harder than model development. Seamless workflow integration drove adoption.

    \item \textbf{Start as Assistant, Not Replacement:} Positioned as "AI-assisted" rather than "AI diagnosis". Radiologists remained in control, reducing resistance.

    \item \textbf{Uncertainty Quantification Essential:} Radiologists valued confidence intervals. High uncertainty prompted careful manual review.

    \item \textbf{Privacy Cannot Be Afterthought:} HIPAA compliance, de-identification, and audit logs required from day one. Added 40\% to development time but non-negotiable.

    \item \textbf{Multi-Site Validation Crucial:} Performance varied across hospitals due to equipment differences. Required site-specific calibration.
\end{enumerate}

\section{Case Study 4: Autonomous Vehicle Perception Stack}

\subsection{Business Context and Requirements}

\textbf{Company:} Autonomous delivery vehicle startup\\
\textbf{Challenge:} Safe navigation in complex urban environments\\
\textbf{Goal:} Level 4 autonomy for last-mile delivery

\textbf{Key Requirements:}

\begin{itemize}
    \item \textbf{Safety:} Zero critical failures, 99.999\% reliability
    \item \textbf{Real-Time:} 100ms perception-to-action latency
    \item \textbf{All-Weather:} Function in rain, fog, night conditions
    \item \textbf{Edge Deployment:} Run on vehicle hardware (limited compute)
    \item \textbf{Continuous Learning:} Improve from fleet data
    \item \textbf{Verification:} Provable safety guarantees
\end{itemize}

\subsection{Perception Pipeline}

\begin{lstlisting}[language=Python, caption=Multi-Sensor Perception System]
#!/usr/bin/env python3
"""
Autonomous Vehicle Perception Stack
Multi-sensor fusion for object detection and tracking
"""

import torch
import torch.nn as nn
import numpy as np
from dataclasses import dataclass
from typing import List, Dict, Tuple, Optional
from enum import Enum
import cv2

class ObjectClass(Enum):
    """Detected object classes"""
    CAR = 0
    PEDESTRIAN = 1
    CYCLIST = 2
    TRUCK = 3
    MOTORCYCLE = 4
    BUS = 5
    TRAFFIC_LIGHT = 6
    STOP_SIGN = 7
    UNKNOWN = 9

@dataclass
class Detection3D:
    """3D object detection"""
    class_id: ObjectClass
    confidence: float
    bbox_3d: np.ndarray  # [x, y, z, length, width, height, yaw]
    velocity: np.ndarray  # [vx, vy, vz]
    tracking_id: Optional[int] = None

@dataclass
class SensorData:
    """Multi-sensor input"""
    camera_images: List[np.ndarray]  # Multiple cameras
    lidar_point_cloud: np.ndarray
    radar_detections: List[Dict]
    imu_data: Dict
    gps_data: Dict
    timestamp: float

class MultiSensorFusionNetwork(nn.Module):
    """Neural network for sensor fusion"""

    def __init__(self):
        super().__init__()

        # Camera branch (per camera)
        self.camera_encoder = self._build_camera_encoder()

        # LiDAR branch
        self.lidar_encoder = self._build_lidar_encoder()

        # Fusion layers
        self.fusion_network = nn.Sequential(
            nn.Linear(512 + 256, 512),
            nn.ReLU(),
            nn.Dropout(0.1),
            nn.Linear(512, 256)
        )

        # Detection head
        self.detection_head = nn.Sequential(
            nn.Linear(256, 128),
            nn.ReLU(),
            nn.Linear(128, 7 + 10)  # 7 bbox params + 10 classes
        )

    def _build_camera_encoder(self) -> nn.Module:
        """Build camera feature encoder"""
        return nn.Sequential(
            # Simplified - in production use ResNet/EfficientNet
            nn.Conv2d(3, 64, kernel_size=7, stride=2, padding=3),
            nn.BatchNorm2d(64),
            nn.ReLU(),
            nn.MaxPool2d(3, stride=2, padding=1),

            nn.Conv2d(64, 128, kernel_size=3, padding=1),
            nn.BatchNorm2d(128),
            nn.ReLU(),
            nn.MaxPool2d(2),

            nn.Conv2d(128, 256, kernel_size=3, padding=1),
            nn.BatchNorm2d(256),
            nn.ReLU(),
            nn.AdaptiveAvgPool2d((1, 1)),
            nn.Flatten(),
            nn.Linear(256, 512)
        )

    def _build_lidar_encoder(self) -> nn.Module:
        """Build LiDAR point cloud encoder"""
        return nn.Sequential(
            # PointNet-style architecture
            nn.Linear(4, 64),  # xyz + intensity
            nn.ReLU(),
            nn.Linear(64, 128),
            nn.ReLU(),
            nn.Linear(128, 256),
            nn.ReLU(),
            # Global max pooling handled separately
        )

    def forward(
        self,
        camera_images: torch.Tensor,
        lidar_points: torch.Tensor
    ) -> torch.Tensor:
        """Forward pass"""

        # Process camera images
        camera_features = []
        for i in range(camera_images.shape[1]):  # Multiple cameras
            img = camera_images[:, i]
            feat = self.camera_encoder(img)
            camera_features.append(feat)

        camera_features = torch.stack(camera_features, dim=1).mean(dim=1)

        # Process LiDAR point cloud
        lidar_features = self.lidar_encoder(lidar_points)
        lidar_features = lidar_features.max(dim=1)[0]  # Global max pool

        # Fuse features
        fused = torch.cat([camera_features, lidar_features], dim=1)
        fused_features = self.fusion_network(fused)

        # Detect objects
        detections = self.detection_head(fused_features)

        return detections

class PerceptionSystem:
    """Complete perception system"""

    def __init__(self, model_path: str, device: str = 'cuda'):
        self.device = device

        # Load detection model
        self.model = MultiSensorFusionNetwork()
        self.model.load_state_dict(torch.load(model_path))
        self.model.to(device)
        self.model.eval()

        # Object tracker
        self.tracker = MultiObjectTracker()

        # Safety validator
        self.safety_validator = SafetyValidator()

    def process_sensor_data(
        self,
        sensor_data: SensorData
    ) -> List[Detection3D]:
        """Process multi-sensor data and detect objects"""

        # Preprocess camera images
        camera_tensors = self._preprocess_cameras(sensor_data.camera_images)

        # Preprocess LiDAR
        lidar_tensor = self._preprocess_lidar(sensor_data.lidar_point_cloud)

        # Run detection network
        with torch.no_grad():
            detections = self.model(camera_tensors, lidar_tensor)

        # Parse detections
        detection_list = self._parse_detections(detections)

        # Fuse with radar
        detection_list = self._fuse_radar(detection_list, sensor_data.radar_detections)

        # Track objects
        tracked_detections = self.tracker.update(detection_list, sensor_data.timestamp)

        # Validate safety
        validated_detections = self.safety_validator.validate(tracked_detections)

        return validated_detections

    def _preprocess_cameras(self, images: List[np.ndarray]) -> torch.Tensor:
        """Preprocess camera images"""

        processed = []
        for img in images:
            # Resize and normalize
            img = cv2.resize(img, (640, 480))
            img = torch.FloatTensor(img).permute(2, 0, 1) / 255.0
            processed.append(img)

        return torch.stack(processed).unsqueeze(0).to(self.device)

    def _preprocess_lidar(self, point_cloud: np.ndarray) -> torch.Tensor:
        """Preprocess LiDAR point cloud"""

        # Filter points (remove ground plane, distant points)
        mask = (point_cloud[:, 2] > -1.5) & (np.linalg.norm(point_cloud[:, :3], axis=1) < 70)
        filtered_points = point_cloud[mask]

        # Subsample to fixed size
        if len(filtered_points) > 16384:
            indices = np.random.choice(len(filtered_points), 16384, replace=False)
            filtered_points = filtered_points[indices]

        # Convert to tensor
        points_tensor = torch.FloatTensor(filtered_points).unsqueeze(0).to(self.device)

        return points_tensor

    def _parse_detections(self, detections: torch.Tensor) -> List[Detection3D]:
        """Parse network output into detections"""

        detections_list = []

        # Simplified - in production, use NMS and proper post-processing
        for i in range(detections.shape[0]):
            det = detections[i].cpu().numpy()

            # Parse bbox and class
            bbox_3d = det[:7]
            class_probs = det[7:]

            class_id = class_probs.argmax()
            confidence = class_probs[class_id]

            if confidence > 0.5:  # Confidence threshold
                detection = Detection3D(
                    class_id=ObjectClass(class_id),
                    confidence=float(confidence),
                    bbox_3d=bbox_3d,
                    velocity=np.zeros(3)  # Will be estimated by tracker
                )
                detections_list.append(detection)

        return detections_list

    def _fuse_radar(
        self,
        vision_detections: List[Detection3D],
        radar_detections: List[Dict]
    ) -> List[Detection3D]:
        """Fuse radar velocity measurements"""

        # Associate radar with vision detections
        for radar_det in radar_detections:
            radar_pos = np.array([radar_det['x'], radar_det['y'], 0])

            # Find closest vision detection
            min_dist = float('inf')
            closest_det = None

            for det in vision_detections:
                vision_pos = det.bbox_3d[:3]
                dist = np.linalg.norm(vision_pos - radar_pos)

                if dist < min_dist:
                    min_dist = dist
                    closest_det = det

            # If close enough, update velocity
            if closest_det and min_dist < 2.0:
                closest_det.velocity[:2] = [radar_det['vx'], radar_det['vy']]

        return vision_detections

class MultiObjectTracker:
    """Track objects across frames"""

    def __init__(self):
        self.tracks: Dict[int, Track] = {}
        self.next_id = 0

    def update(
        self,
        detections: List[Detection3D],
        timestamp: float
    ) -> List[Detection3D]:
        """Update tracks with new detections"""

        # Associate detections with existing tracks
        assignments, unassigned_detections = self._associate(detections)

        # Update existing tracks
        for track_id, detection in assignments.items():
            self.tracks[track_id].update(detection, timestamp)
            detection.tracking_id = track_id

        # Create new tracks
        for det_idx in unassigned_detections:
            new_track = Track(self.next_id, detections[det_idx], timestamp)
            self.tracks[self.next_id] = new_track
            detections[det_idx].tracking_id = self.next_id
            self.next_id += 1

        # Remove stale tracks
        self._remove_stale_tracks(timestamp)

        return detections

    def _associate(
        self,
        detections: List[Detection3D]
    ) -> Tuple[Dict[int, Detection3D], List[int]]:
        """Associate detections with tracks using Hungarian algorithm"""

        # Simplified - use IoU matching
        # In production, use Kalman filter predictions and Hungarian algorithm

        assignments = {}
        unassigned = list(range(len(detections)))

        for track_id, track in self.tracks.items():
            best_iou = 0
            best_det_idx = None

            for det_idx, detection in enumerate(detections):
                iou = self._calculate_iou_3d(track.bbox_3d, detection.bbox_3d)

                if iou > best_iou and iou > 0.3:
                    best_iou = iou
                    best_det_idx = det_idx

            if best_det_idx is not None:
                assignments[track_id] = detections[best_det_idx]
                if best_det_idx in unassigned:
                    unassigned.remove(best_det_idx)

        return assignments, unassigned

    def _calculate_iou_3d(self, bbox1: np.ndarray, bbox2: np.ndarray) -> float:
        """Calculate 3D IoU (simplified bird's eye view)"""

        # Extract x, y, length, width
        x1, y1, l1, w1 = bbox1[0], bbox1[1], bbox1[3], bbox1[4]
        x2, y2, l2, w2 = bbox2[0], bbox2[1], bbox2[3], bbox2[4]

        # Calculate intersection
        x_overlap = max(0, min(x1 + l1/2, x2 + l2/2) - max(x1 - l1/2, x2 - l2/2))
        y_overlap = max(0, min(y1 + w1/2, y2 + w2/2) - max(y1 - w1/2, y2 - w2/2))

        intersection = x_overlap * y_overlap

        # Calculate union
        area1 = l1 * w1
        area2 = l2 * w2
        union = area1 + area2 - intersection

        return intersection / union if union > 0 else 0

    def _remove_stale_tracks(self, current_time: float):
        """Remove tracks not updated recently"""

        stale_tracks = []
        for track_id, track in self.tracks.items():
            if current_time - track.last_update_time > 1.0:  # 1 second
                stale_tracks.append(track_id)

        for track_id in stale_tracks:
            del self.tracks[track_id]

@dataclass
class Track:
    """Individual object track"""
    track_id: int
    bbox_3d: np.ndarray
    velocity: np.ndarray
    last_update_time: float
    detection_count: int = 1

    def __init__(self, track_id: int, detection: Detection3D, timestamp: float):
        self.track_id = track_id
        self.bbox_3d = detection.bbox_3d
        self.velocity = detection.velocity
        self.last_update_time = timestamp
        self.detection_count = 1

    def update(self, detection: Detection3D, timestamp: float):
        """Update track with new detection"""

        # Smooth bbox (exponential moving average)
        alpha = 0.7
        self.bbox_3d = alpha * detection.bbox_3d + (1 - alpha) * self.bbox_3d

        # Update velocity
        dt = timestamp - self.last_update_time
        if dt > 0:
            velocity_estimate = (detection.bbox_3d[:3] - self.bbox_3d[:3]) / dt
            self.velocity = alpha * velocity_estimate + (1 - alpha) * self.velocity

        self.last_update_time = timestamp
        self.detection_count += 1

class SafetyValidator:
    """Validate detections for safety-critical applications"""

    def __init__(self):
        self.min_confidence = 0.7
        self.max_reasonable_velocity = 50.0  # m/s

    def validate(self, detections: List[Detection3D]) -> List[Detection3D]:
        """Validate and filter detections"""

        validated = []

        for detection in detections:
            # Check confidence
            if detection.confidence < self.min_confidence:
                continue

            # Check velocity is reasonable
            velocity_magnitude = np.linalg.norm(detection.velocity)
            if velocity_magnitude > self.max_reasonable_velocity:
                continue

            # Check bbox is within sensor range
            distance = np.linalg.norm(detection.bbox_3d[:3])
            if distance > 100:  # 100 meters
                continue

            validated.append(detection)

        return validated

# Example usage
if __name__ == "__main__":
    perception = PerceptionSystem(
        model_path='models/perception_fusion.pth',
        device='cuda'
    )

    # Simulate sensor data
    sensor_data = SensorData(
        camera_images=[np.zeros((480, 640, 3)) for _ in range(4)],
        lidar_point_cloud=np.random.randn(50000, 4),
        radar_detections=[],
        imu_data={},
        gps_data={},
        timestamp=0.0
    )

    # Process
    detections = perception.process_sensor_data(sensor_data)

    print(f"Detected {len(detections)} objects")
    for det in detections:
        print(f"  {det.class_id.name}: confidence={det.confidence:.2f}, "
              f"position=({det.bbox_3d[0]:.1f}, {det.bbox_3d[1]:.1f})")
\end{lstlisting}

\subsection{Results and Impact}

\textbf{Performance Metrics:}

\begin{itemize}
    \item \textbf{Detection Accuracy:} 98.2\% recall, 96.8\% precision
    \item \textbf{Latency:} 78ms end-to-end (perception to planning)
    \item \textbf{Safety:} Zero critical failures in 2M test miles
    \item \textbf{All-Weather:} 94\% performance maintained in adverse conditions
    \item \textbf{Intervention Rate:} 0.12 per 1000 miles (industry-leading)
\end{itemize}

\textbf{Business Impact:}

\begin{itemize}
    \item \textbf{Deployment:} 500 vehicles operational in 5 cities
    \item \textbf{Deliveries:} 50,000 autonomous deliveries per week
    \item \textbf{Cost Savings:} 60\% reduction vs. human drivers
    \item \textbf{Efficiency:} 24/7 operation, 40\% more deliveries per vehicle
\end{itemize}

\subsection{Lessons Learned}

\begin{enumerate}
    \item \textbf{Sensor Fusion is Non-Negotiable:} No single sensor modality is sufficient. Camera-LiDAR-radar fusion caught 18\% more edge cases than camera-only.

    \item \textbf{Edge Cases Dominate Development:} 90\% of engineering time on <1\% of scenarios. Construction zones, emergency vehicles, and edge weather conditions required specific handling.

    \item \textbf{Simulation Accelerated Development:} Synthetic data and simulation generated 10M training examples. Real-world data collection alone would take years.

    \item \textbf{Continuous Learning Essential:} Fleet learning improved performance 15\% over 6 months. Automated pipeline to retrain from fleet data weekly.

    \item \textbf{Safety Requires Redundancy:} Dual perception systems with different architectures. Discrepancies trigger safe stop.

    \item \textbf{Latency Matters More Than Accuracy:} 98\% accurate in 100ms better than 99\% in 200ms for real-time control.

    \item \textbf{Hardware-Software Co-Design:} Custom accelerators reduced latency 3x and power consumption 5x vs. off-the-shelf GPUs.

    \item \textbf{Verification is Harder Than Development:} Proving safety to regulators required formal verification methods and extensive testing.
\end{enumerate}

\section{Cross-Cutting Themes and Best Practices}

Across all four case studies, several common patterns emerged:

\textbf{Universal Best Practices:}

\begin{enumerate}
    \item \textbf{Start Simple, Then Optimize:} All systems began with simple baselines before adding complexity

    \item \textbf{Measure What Matters:} Offline metrics often misaligned with business objectives. Online A/B testing essential.

    \item \textbf{Monitoring is Not Optional:} Production issues were caught by monitoring, not caught in testing

    \item \textbf{Data Quality Trumps Algorithm Choice:} Better data improved results more than better models

    \item \textbf{Incremental Deployment Reduces Risk:} Gradual rollouts caught issues before full deployment

    \item \textbf{Team Diversity Matters:} Cross-functional teams (ML, engineering, domain experts) built better systems

    \item \textbf{Documentation Pays Off:} Well-documented systems were easier to maintain and debug

    \item \textbf{Plan for Failure:} Systems with graceful degradation and fallbacks were more reliable
\end{enumerate}

\section{Chapter Summary}

This chapter presented four comprehensive industry case studies demonstrating end-to-end ML system implementations across diverse domains. Key takeaways:

\begin{itemize}
    \item \textbf{Domain Adaptation:} ML patterns adapt to specific industry requirements
    \item \textbf{Real-World Complexity:} Production systems face challenges beyond model accuracy
    \item \textbf{Hybrid Approaches:} Combining multiple techniques often outperforms single solutions
    \item \textbf{Continuous Improvement:} All successful systems evolved through iteration
    \item \textbf{Cross-Functional Success:} Technical excellence alone is insufficient without domain expertise and business alignment
\end{itemize}

\section{Exercises}

\begin{exercise}
\textbf{E-Commerce Recommendation System Analysis}

Analyze the recommendation system case study and:

\begin{enumerate}
    \item Calculate the expected ROI if deployed at a company with 2M monthly active users and \$35 average order value
    \item Design an A/B test to measure the impact of adding diversity to recommendations
    \item Propose a solution for the cold start problem for new users with zero history
    \item Design a feature for handling seasonal trends (holidays, sales events)
\end{enumerate}

\textbf{Deliverables:}
\begin{itemize}
    \item ROI calculation with assumptions documented
    \item A/B test design including sample size, duration, and metrics
    \item Cold start strategy with pseudocode
    \item Seasonal adaptation architecture diagram
\end{itemize}
\end{exercise}

\begin{exercise}
\textbf{Fraud Detection System Design}

Design fraud detection improvements:

\begin{enumerate}
    \item Identify three new features not mentioned in the case study that could improve fraud detection
    \item Design a strategy to handle concept drift (fraud patterns changing over time)
    \item Propose a method to reduce false positives for high-value customers
    \item Design an active learning system to efficiently label uncertain transactions
\end{enumerate}

\textbf{Deliverables:}
\begin{itemize}
    \item Feature specifications with rationale
    \item Concept drift detection and adaptation strategy
    \item VIP customer handling workflow
    \item Active learning pipeline design
\end{itemize}
\end{exercise}

\begin{exercise}
\textbf{Healthcare System Compliance Analysis}

Evaluate the healthcare diagnostic system:

\begin{enumerate}
    \item Research FDA 510(k) clearance requirements and outline steps for submission
    \item Design a clinical validation study to demonstrate safety and efficacy
    \item Propose metrics beyond accuracy that matter for radiologist adoption
    \item Design a monitoring system to detect model degradation in production
\end{enumerate}

\textbf{Deliverables:}
\begin{itemize}
    \item FDA submission outline with timeline
    \item Clinical validation protocol
    \item User experience metrics framework
    \item Model monitoring dashboard specification
\end{itemize}
\end{exercise}

\begin{exercise}
\textbf{Autonomous Vehicle Safety Analysis}

Analyze the perception system safety:

\begin{enumerate}
    \item Identify potential failure modes and their consequences
    \item Design a safety validation test suite with coverage goals
    \item Propose a simulation strategy to test edge cases
    \item Design redundancy mechanisms for critical components
\end{enumerate}

\textbf{Deliverables:}
\begin{itemize}
    \item FMEA (Failure Mode and Effects Analysis) table
    \item Test suite specification with success criteria
    \item Simulation architecture and scenario generation strategy
    \item System redundancy design document
\end{itemize}
\end{exercise}
